% Generated draft: requires mathematical and source review.
\documentclass[11pt]{article}
\usepackage[margin=25mm]{geometry}
\usepackage{fontspec}
\setmainfont{Latin Modern Roman}
\usepackage{amsmath,amssymb,mathtools,unicode-math}
\setmathfont{Latin Modern Math}
\usepackage{enumitem,needspace,xurl,hyperref}
\setlength{\parindent}{0pt}
\setlength{\parskip}{5pt}
\setlength{\emergencystretch}{3em}
\newcommand{\N}{\mathbb{N}}
\newcommand{\Z}{\mathbb{Z}}
\newcommand{\Q}{\mathbb{Q}}
\newcommand{\R}{\mathbb{R}}
\newcommand{\C}{\mathbb{C}}
\newcommand{\F}{\mathbb{F}}
\newcommand{\A}{\mathbb{A}}
\newcommand{\PP}{\mathbb{P}}
\newcommand{\E}{\mathbb{E}}
\newcommand{\eps}{\varepsilon}
\newcommand{\dd}{\,\mathrm{d}}
\begin{document}
\section*{UnsolvedMath ID: 20003192}
\textbf{Source problem code:} catalogue entry
\subsection{Definitions and mathematical statement}
Marc Burger asks to count rational representations whose images lie in the integer points of the symplectic group. The canonical sentence gives no domain group, symplectic rank, rationality convention, equivalence relation, height or asymptotic parameter. Nearby material places it in the surface-group setting, but this does not uniquely determine a counting problem. In particular, merely having integral matrices already implies rational matrix entries; ``rational representation'' could have other meanings requiring clarification.

The aid fixes integers $g,m\geq1$, a closed oriented genus-$g$ surface and the group
\[
 \pi_1(\Sigma_g)=\langle a_1,b_1,\ldots,a_g,b_g:
 \prod_{j=1}^g[a_j,b_j]=1\rangle,
\]
with target $Sp_{2m}(\Z)$ preserving a fixed standard integral symplectic form. In this interpretation the literal set is
\[
 \operatorname{Hom}(\pi_1(\Sigma_g),Sp_{2m}(\Z)).
\]
The supplied result says it is countably infinite for every such $g,m$, not a finite count. It also reports a primitive cyclic-image family whose members remain inequivalent even after real symplectic conjugacy and mapping-class precomposition, by trace separation of their image generators. The actual matrices and parameterization of that family are absent from the record. Primitive here must be checked in the construction and is not silently interpreted as nondivisibility in the ambient arithmetic group.

A separate, explicitly chosen finite-count reformulation fixes the displayed generators and defines
\[
 H(\rho)=\max\{|\rho(a_j)_{uv}|,|\rho(b_j)_{uv}|:
 1\leq j\leq g,\ 1\leq u,v\leq2m\}.
\]
Let $[\rho]$ be its orbit under integral target conjugacy and mapping-class precomposition, with an orientation convention for the mapping class group fixed in advance. Set
\[
 h([\rho])=\min_{\rho'\in[\rho]}H(\rho'),\qquad
 N_{g,m}(B)=\#\{[\rho]:h([\rho])\leq B\}.
\]
The minimum is over the algebraic orbit and is achieved as a minimum of nonnegative integer heights; no effective orbit-minimization algorithm is asserted. The aid reports, for every integer $B\geq3$,
\[
 B-2\leq N_{g,m}(B)\leq(2B+1)^{8gm^2}.
\]
There are $8gm^2$ entries in the $2g$ generator matrices, explaining the crude finite upper bound before imposing symplectic and surface relations. The lower bound is attributed to the trace-separated cyclic family and is reported to survive the larger real-conjugacy equivalence as well.

This gives finiteness at bounded minimal generator height and coarse quantitative bounds in one convention. It supplies neither an exact count nor a sharp growth asymptotic, and it is not claimed to be the uniquely intended meaning of the original question. Genus zero, rank zero, different rationality notions and other quotient or height choices are outside the stated result.
\subsection{Short English statement}
The literal integral-representation set in the aid’s fixed genus and rank is countably infinite. A finite counting problem appears only after choosing a height and equivalence relation. The stated bounds concern that particular minimal-generator-height count, not an exact or canonical interpretation of the original request.
\subsection{Sources}
\begin{itemize}
\item Local UnsolvedMath registry, permanent ID 20003192, source code catalogue entry; \texttt{aim-topology-notes.json}, notes[279], source tag question.
\item AIM, \emph{Representations of surface groups}, Question 5.4, Marc Burger: \url{https://aimath.org/WWN/surfacegroups/surfacegroups.pdf}. The supplied verification reports checking the PDF and using nearby questions for the domain context; this draft does not independently repeat it.
\item The supplied research summary gives the cyclic family and height bounds without an independent authoritative external bibliography or explicit family matrices.
\item Machine-generated research attempt 1: partial result, original status unknown at the run, rigor labelled ``proved,'' low novelty confidence; registry partially solved. Not peer reviewed; independently verify before citation. Registry updated 14 August 2026.
\end{itemize}
\end{document}
